\section{Sumas de más de dos términos}

Hasta ahora hemos definido la adición como una operación binaria:
\[
+\colon\mathbb{Z}\times\mathbb{Z}\to\mathbb{Z}.
\]
Esto significa que la operación \(+\) recibe exactamente dos números cada vez que se aplica. Por ejemplo,
\[
3+5
\]
es una aplicación de la adición a los números \(3\) y \(5\).
Sin embargo, es habitual encontrar expresiones como
\[
3+5+7
\]
o incluso
\[
12+4+(-8)+15+(-2).
\]

A primera vista podría parecer que ahora estamos sumando tres, cuatro o cinco números simultáneamente. No es así. La adición continúa siendo una operación binaria: lo que hacemos es aplicarla varias veces.

\subsection{Sumar de a dos}

Considere la expresión
\[
3+5+7.
\]
Una manera de interpretarla es realizar primero
\[
3+5
\]
y después sumar \(7\) al número obtenido:
\[
(3+5)+7.
\]
Primero tenemos
\[
3+5=8,
\]
y luego
\[
8+7=15.
\]
Por lo tanto,
\[
(3+5)+7=15.
\]
También podríamos comenzar sumando \(5\) y \(7\):
\[
3+(5+7).
\]
En este caso,
\[
5+7=12
\]
y luego
\[
3+12=15.
\]
Obtenemos nuevamente el mismo número:
\[
3+(5+7)=15.
\]

Esto no ocurre por casualidad. La adición posee una propiedad llamada \emph{asociatividad}.

\subsection{La propiedad asociativa}

Para cualesquiera números enteros \(a\), \(b\) y \(c\),
\[
(a+b)+c=a+(b+c).
\]
Esta propiedad nos dice que, cuando únicamente estamos sumando, podemos cambiar la manera en que agrupamos los sumandos sin modificar el número representado.
Por ejemplo,
\[
(2+8)+5=2+(8+5).
\]
En ambos casos obtenemos
\[
15.
\]
Gracias a esta propiedad, normalmente omitimos los paréntesis y escribimos simplemente
\[
a+b+c.
\]
Por convención, esta escritura puede entenderse como aplicaciones sucesivas de la operación de adición. La propiedad asociativa garantiza que la forma de agruparlas no altera su valor.

Por esta razón podemos escribir sin ambigüedad expresiones como
\[
2+5+7+9.
\]

La operación sigue aplicándose de a dos números cada vez; simplemente hemos evitado escribir todos los paréntesis que serían necesarios para mostrar cada aplicación individual.

\subsection{Sumandos y términos}

En una expresión como
\[
2+5+7+9,
\]
los números
\[
2,\qquad5,\qquad7,\qquad9
\]
son los \emph{sumandos} de la adición.

También utilizaremos frecuentemente la palabra \emph{término} para referirnos a cada una de las partes principales que componen una expresión.

En una suma como la anterior, cada término es también un sumando. Más adelante encontraremos expresiones en las que un término puede ser algo más complejo que un único número.
Por ejemplo, en
\[
12+(3\cdot 4)+7,
\]
podemos considerar como términos principales
\[
12,\qquad 3\cdot4,\qquad7.
\]
Por ahora trabajaremos principalmente con términos que son números enteros.

\subsection{Expresiones con sustracciones}

La situación requiere un poco más de atención cuando aparecen signos de sustracción.
Considere
\[
123+3-2+(-20).
\]
Sabemos que toda sustracción puede escribirse como la suma del opuesto. Por lo tanto,
\[
123+3-2+(-20)
\]
puede reescribirse como
\[
123+3+(-2)+(-20).
\]
Ahora la expresión contiene únicamente adiciones.
Sus sumandos son
\[
123,\qquad3,\qquad-2,\qquad-20.
\]
Observe especialmente que uno de los sumandos es
\[
-2,
\]
no \(2\).

Una vez escrita toda la expresión como una suma, podemos aprovechar las propiedades de la adición para agrupar y ordenar los sumandos de la manera que resulte más conveniente.
Por ejemplo,
\[
123+3+(-2)+(-20)
\]
puede reorganizarse como
\[
123+(-20)+3+(-2).
\]
Luego podemos agrupar:
\[
\bigl(123+(-20)\bigr)+\bigl(3+(-2)\bigr),
\]
y evaluar:
\[
103+1=104.
\]
Por lo tanto,
\[
123+3-2+(-20)=104.
\]

\subsection{El signo forma parte del sumando}

Esta forma de interpretar una expresión permite aclarar algo que será muy importante durante los cálculos.
En
\[
7-4+9-12,
\]
podemos reemplazar las sustracciones por sumas de opuestos:
\[
7+(-4)+9+(-12).
\]
Ahora podemos identificar claramente los cuatro sumandos:
\[
7,\qquad-4,\qquad9,\qquad-12.
\]
Al reorganizar la suma, cada número conserva su signo:
\[
7+9+(-4)+(-12).
\]
Podemos incluso agrupar los positivos y los negativos:
\[
(7+9)+\bigl((-4)+(-12)\bigr).
\]
Así obtenemos
\[
16+(-16)=0.
\]
Esta manera de trabajar evita pensar que los signos se desplazan independientemente de los números. En una suma de enteros, números como
\[
-4
\qquad\text{y}\qquad
-12
\]
son sumandos completos.

\subsection{Una precaución con la sustracción}

La libertad para reagrupar aparece una vez que hemos expresado todo mediante adiciones.

La sustracción, por sí misma, no es asociativa.
Por ejemplo,
\[
(10-3)-2=7-2=5,
\]
mientras que
\[
10-(3-2)=10-1=9.
\]
Por lo tanto,
\[
(10-3)-2\neq10-(3-2).
\]
Cuando escribimos
\[
10-3-2,
\]
la convención habitual es interpretar las operaciones de izquierda a derecha:
\[
(10-3)-2.
\]
Sin embargo, podemos evitar esta dificultad reescribiendo primero las sustracciones:
\[
10-3-2
=
10+(-3)+(-2).
\]
Ahora sí tenemos únicamente una suma y podemos aplicar libremente las propiedades asociativa y conmutativa:
\[
10+(-3)+(-2)
=
10+\bigl((-3)+(-2)\bigr)
=
10+(-5)
=
5.
\]

\begin{note}
No debe confundirse la posibilidad de reorganizar una suma con la posibilidad de reorganizar arbitrariamente una expresión que contiene sustracciones.

Primero podemos transformar
\[
a-b+c-d
\]
en
\[
a+(-b)+c+(-d).
\]
Una vez hecho esto, tenemos una suma cuyos sumandos son
\[
a,\qquad -b,\qquad c,\qquad -d,
\]
y entonces sí podemos agruparlos u ordenarlos utilizando las propiedades de la adición.
\end{note}

\subsection{Una forma práctica de leer estas expresiones}

Ante una expresión como
\[
25-8+3-(-4)-10,
\]
puede resultar útil realizar primero una tarea puramente simbólica: identificar qué números se están sumando.

Reemplazamos cada sustracción por la suma del opuesto:
\[
25+(-8)+3+\bigl(-(-4)\bigr)+(-10).
\]
Como
\[
-(-4)=4,
\]
obtenemos
\[
25+(-8)+3+4+(-10).
\]
Ahora podemos leer la expresión como una suma de cinco términos:
\[
25,\qquad-8,\qquad3,\qquad4,\qquad-10.
\]
A partir de aquí podemos escoger una agrupación conveniente para evaluarla:
\[
25+(-10)+(-8)+3+4.
\]
Por ejemplo,
\[
\bigl(25+(-10)\bigr)+\bigl((-8)+3\bigr)+4
=
15+(-5)+4
=
14.
\]

El procedimiento puede resumirse de esta manera:
\begin{enumerate}
\item Reescriba cada sustracción como la suma de un opuesto.
\item Identifique los sumandos de la expresión.
\item Simplifique los dobles opuestos cuando aparezcan.
\item Agrupe u ordene los sumandos de una manera conveniente.
\item Evalúe las adiciones sucesivamente.
\end{enumerate}

No estamos introduciendo una nueva forma de sumar. Seguimos utilizando la misma operación binaria que ya conocemos, aplicada repetidas veces.
